\subsection{Implantação}
\label{subsec:implantacao}

A implantação do Movy reflete a separação em camadas descrita na Seção~\ref{sec:arquitetura}: cada artefato é publicado em uma plataforma especializada no seu papel: cliente e API comunicam-se por HTTP/JSON, e a API acessa o banco por meio de uma conexão PostgreSQL sobre SSL. A topologia de execução encadeia, do usuário ao dado, quatro elementos: o navegador acessa o \textit{Worker} da Cloudflare (cliente em execução de borda) por HTTP/JSON; este consome o \textit{Web Service} do Render (API NestJS, servidor persistente), também por HTTP/JSON; e a API conecta-se ao PostgreSQL gerenciado no Neon por TCP/SSL com \textit{pooler}. O ambiente local de desenvolvimento, no qual o \textit{Docker Compose} sobe, em conjunto, o PostgreSQL e a API, permanece inalterado e passa a coexistir com esse ambiente online.\footnote{Implantação validada nas URLs públicas \texttt{movy-app.rn-labs.workers.dev} (cliente) e \texttt{movy-api-av8q.onrender.com} (API).} A Tabela~\ref{tab:implantacao} resume o papel de cada alvo de implantação.

\begin{table}[H]
\centering
\caption{Alvos de implantação do sistema Movy}
\label{tab:implantacao}
\begin{tabularx}{\textwidth}{l l X}
\toprule
\textbf{Camada} & \textbf{Plataforma} & \textbf{Papel na implantação} \\
\midrule
Aplicação cliente & Cloudflare Workers & Hospedagem do \textit{frontend} TanStack Start em execução de borda, próxima ao usuário. \\
API REST & Render (\textit{Web Service}) & Execução contínua do servidor NestJS, recebendo requisições HTTP públicas. \\
Banco de dados & Neon (PostgreSQL gerenciado) & Persistência relacional, acessada exclusivamente pela API. \\
\bottomrule
\end{tabularx}
\end{table}

\subsubsection{Backend, API no Render e banco no Neon}

Para o banco de dados, optou-se pelo Neon, serviço de PostgreSQL gerenciado, em detrimento de alternativas como o Supabase. Embora ambos forneçam PostgreSQL, o Supabase agrega serviços acoplados, autenticação, APIs geradas e \textit{storage}, que se sobreporiam a responsabilidades já implementadas no servidor. A escolha do Neon, focado no banco, preserva o \textit{backend} NestJS como único responsável por autenticação, autorização, regras de negócio e fronteiras de \textit{tenant}, mantendo a coerência com a \textit{Clean Architecture} e o DDD adotados (Seção~\ref{sec:arquitetura}). O Neon expõe duas \textit{strings} de conexão: uma direta, empregada pelos comandos do Prisma CLI na aplicação das \textit{migrations}, e uma com \textit{pooler}, usada pela aplicação em \textit{runtime} para suportar conexões concorrentes. O banco online foi criado vazio e recebeu o histórico completo do \textit{schema} pelo mesmo mecanismo de \textit{migrations} versionadas e \textit{seed} descrito na Seção~\ref{subsec:migracoes}, sem migração de dados a partir do ambiente local, o que é adequado ao objetivo de validar fluxos em ambiente online, partindo de uma base consistente com o \textit{schema} atual.

A API foi publicada no Render, na modalidade \textit{Web Service}, adequada a um servidor que precisa receber requisições HTTP públicas de forma contínua. A opção pelo Render, em vez da Vercel, decorre da natureza da aplicação: o \textit{backend} é um servidor NestJS tradicional, com processo HTTP persistente, acesso ao banco via Prisma e \textit{jobs} agendados internos (\texttt{@nestjs/schedule}), um modelo distinto do de \textit{functions} efêmeras para o qual a Vercel é mais adequada. O arquivo \texttt{docker-compose.yml} não é utilizado no ambiente online; ele permanece restrito ao desenvolvimento local, onde sobe API e banco em conjunto. Em produção, o banco é externo (Neon) e a API executa isolada no Render, que injeta a porta de escuta por variável de ambiente, atendida pela aplicação ao escutar em \texttt{process.env.PORT} e no \textit{host} \texttt{0.0.0.0}.

A passagem do ambiente local conteinerizado para o Render exigiu ajustes para tornar o \textit{startup} previsível em produção. Como o diretório de artefatos gerados pelo Prisma não é versionado, o \textit{script} de \textit{build} passou a gerar o Prisma Client explicitamente antes da compilação (\texttt{prisma generate \&\& nest build}). O comando de \textit{start}, por sua vez, foi alterado de \texttt{nest start}, adequado ao desenvolvimento, porém pesado para \textit{runtime}, para a execução direta do JavaScript compilado (\texttt{node dist/src/main.js}), após o primeiro \textit{deploy} apresentar esgotamento de memória (\textit{JavaScript heap out of memory}) em instância pequena. Por fim, ajustes de exposição reduziram a superfície pública: a documentação Swagger passou a ser registrada apenas quando habilitada explicitamente pela variável \texttt{ENABLE\_SWAGGER} (mantida desligada em produção, para não facilitar a descoberta das rotas), e a rota raiz, que antes retornava uma mensagem de desenvolvimento, passou a responder um \textit{health check} discreto.

A configuração do serviço no Render segue o mesmo contrato de variáveis de ambiente da aplicação: a \texttt{DATABASE\_URL} (apontando para a conexão \textit{pooled} do Neon), o segredo de assinatura dos \textit{tokens} (\texttt{JWT\_SECRET}, mantido apenas como variável secreta e nunca versionado), os tempos de expiração de \textit{access} e \textit{refresh tokens}, a lista de origens autorizadas para CORS (\texttt{ALLOWED\_ORIGINS}) e o interruptor do Swagger. Em termos de segurança operacional, registra-se que o CORS não é tratado como barreira principal, a autorização efetiva permanece a cargo dos \textit{guards} do servidor (\texttt{JwtAuthGuard}, \texttt{RolesGuard}, \texttt{TenantFilterGuard} e \texttt{DevGuard}), e que a correção automática de vulnerabilidades (\texttt{npm audit fix}) foi deliberadamente evitada no \textit{deploy}, por alterar o \textit{lockfile} e poder introduzir \textit{breaking changes}, comprometendo a reprodutibilidade; o tratamento dessas vulnerabilidades foi remetido a uma tarefa específica, com testes.

\subsubsection{Frontend, aplicação cliente na Cloudflare Workers}

A aplicação cliente foi publicada na Cloudflare Workers, alvo de implantação para o qual o projeto já estava configurado: o arquivo \texttt{wrangler.jsonc} define o \textit{Worker} e o \textit{preset} de \textit{build} baseado em Vite já inclui o adaptador que gera os artefatos consumidos pela ferramenta de linha de comando Wrangler. A publicação como aplicação Vite estática convencional foi descartada porque o \textit{build} do TanStack Start não produz uma saída estática única (um \texttt{index.html} na raiz), mas artefatos separados de cliente e de servidor (\texttt{dist/client/} e \texttt{dist/server/}); servi-la como sítio estático resultou em erro \texttt{404}. Manteve-se, assim, o caminho já suportado pelo projeto, evitando adaptações adicionais de \textit{runtime}.

Uma particularidade da camada de apresentação é que a URL base da API (\texttt{VITE\_API\_URL}) é embutida no \textit{bundle} em tempo de \textit{build}, e não lida em \textit{runtime}. Por isso, o valor precisa estar correto antes da geração do \textit{build} de produção: em desenvolvimento, aponta para a API local; na publicação, para a API hospedada no Render. O processo de \textit{deploy} resume-se a configurar essa variável, gerar o \textit{build}, autenticar a CLI Wrangler na conta Cloudflare e publicar o \textit{Worker}, que passa a responder em um subdomínio \texttt{workers.dev} registrado para a conta.

A integração entre as duas camadas publicadas depende de um único ajuste de CORS no servidor: a origem do \textit{Worker} precisa constar em \texttt{ALLOWED\_ORIGINS} na API, sob pena de as chamadas do navegador serem bloqueadas. Após a publicação, a validação seguiu o roteiro de verificação manual (Seção~\ref{subsec:verificacao}) aplicado ao ambiente online: carregamento da página inicial, listagem pública de viagens, \textit{login}, ausência de erros de CORS nas chamadas à API e exercício dos fluxos principais, cadastro de passageiro e de empresa, reserva de viagem, painel administrativo e área do motorista.
